\section{Stores Ahead: Upwards Proved, Downwards Diagnosed}\label{sec:rewnif}

\subsection{The programs}
The test \texttt{forall/rewnif} writes one cell ahead before it writes the current one; \texttt{rewnifrev}
(under \texttt{forall} and \texttt{min}) does the same in a loop that runs downwards. Both assert
\texttt{a[i] >= i} for all cells.
\begin{lstlisting}[language=C]
for (i = 0; i < SIZE; i++) { if (i + 1 < SIZE) a[i+1] = i; a[i] = i; }        // rewnif
for (i = SIZE - 1; i >= 0; i--) { if (i - 1 >= 0) a[i-1] = i - 2; a[i] = i; }  // rewnifrev
for (i = 0; i < SIZE; i++) __VERIFIER_assert(a[i] >= i);
\end{lstlisting}
The runs use the module \texttt{strides} (\texttt{-include teq:strides}), which holds the lemmas on stores
ahead. The summary of the MonoCera tests records \texttt{rewnif} as closed and \texttt{rewnifrev} as the one
test whose lemma ``has the same form as the trace in the goal when printed, yet does not match'', with the
literal $-1$ excluded as the cause. This section re-traces the first and resolves the second.

\subsection{Upwards: the trace equation system and the proof}
With $S$ for \texttt{SIZE}, $1 < S$ from the harness, the filling loop is
\begin{align}
  A &= \mathsf{trace}(a_0, \mathsf{for}\ \bl{K}, \bl{I}\ \mathsf{do}\ a, k, i \mapsto \upd{i + 1 < S}{a[i + 1 := i][i := i]}\ \ \upd{\neg(i + 1 < S)}{a[i := i]}) \label{eq:rw1}\\
  I &= \tr{0}{\bl{K}}{i, k \mapsto \upd{i + 1 < S}{i + 1}\ \ \upd{\neg(i + 1 < S)}{i + 1}} \nonumber\\
  &\al{\bl{K}, \bl{A}, \bl{I}}{k, a, i \mapsto i < S} \nonumber
\end{align}
and the lemma of the module is
\begin{align*}
  &\lemma{select\_current\_trace\_ahead}:\quad \intf(n) \le m\\
  &\quad\Longrightarrow\ \cur(\mathsf{trace}(a_0, \mathsf{for}\ \bl{\cnt(0, 1, n)}, \bl{\cnt(0, 1, n)}\ \mathsf{do}\ a, k, i \mapsto {}\\
  &\qquad\qquad \upd{k + 1 < m}{a[k + 1 := F(k)][k := G(k)]}\ \ \upd{m = k + 1}{a[k := G(k)]}))[j]\\
  &\qquad = \mathsf{if}\ 0 \le j < \intf(n)\ \mathsf{then}\ G(j)\\
  &\qquad\quad \mathsf{else\ if}\ j = \intf(n) \wedge 0 < j < m\ \mathsf{then}\ F(j - 1)\ \mathsf{else}\ a_0[j]
\end{align*}
The value written ahead at $k + 1$ is overwritten by the next iteration, so a cell $j < n$ holds what
iteration $j$ writes last, and the cell $n$ holds what the last iteration wrote ahead, if it did. The side
condition $\intf(n) \le m$ makes the two guards cover all iterations: for $k < n \le m$, either $k + 1 < m$
or $k + 1 = m$. The proof takes 27 rule applications: the counter $I$ is merged and recognized, the
constraint gives $\intf(n) \le S$, which is the side condition, and \lemma{select\_current\_trace\_ahead}
fires twice on the reads $a'[k]$ of the checking loop's constraint and $a'[i'']$ of the conclusion, with
$m := S$ and $F := G := \lambda z.\ z$ (the stored values $i$ abstracted over the bound variable), giving
$\ite{0 \le j < \intf(n)}{j}{\ite{j = \intf(n) \wedge 0 < j < S}{j - 1}{a_0[j]}}$; the exit
$i'' = \intf(n'')$ with $\intf(n'') = S = \intf(n)$ puts the conclusion's index in the first branch
(\lemma{integer\_injective}), and $j \ge j$ holds.

\subsection{Downwards: the goal that remained}
For \texttt{rewnifrev} the counter runs from $S - 1$ down, and the simplifier's goal, after 54 rule
applications and with the module's lemma for the downward loop present but silent, is
\begin{align}
  A &= \mathsf{trace}(a_0, \mathsf{for}\ \bl{K}, \bl{I}\ \mathsf{do}\ a, k, i \mapsto \upd{1 \le i}{a[i + (-1) := i + (-2)][i := i]}\ \ \upd{i < 1}{a[i := i]}) \label{eq:rw2}\\
  a' &= \cur(A) \nonumber\\
  I &= \cnt(\intf(n) + i', -1, n) \qquad 0 = i' + 1 \qquad 0 < \intf(n) + i' < 99999 \label{eq:rw3}\\
  &\mathsf{above\_iterations}(n'', \mapc{\bl{K''}}{k \mapsto a'[k] + 1}) \label{eq:rw4}\\
  &i'' = \intf(n'') \le \intf(n) + i' \nonumber\\
  &\Longrightarrow\ i'' \le a'[i''] \nonumber
\end{align}
Two things have happened to the facts. The exit $\neg(0 \le i')$ with the constraint $0 \le i$ gives
$i' = -1$, and the simplifier has eliminated $S$ in favour of the counter: $S - 1 = \intf(n) + i'$, so the
start of $I$ is the term $\intf(n) + i'$ (with $i' = -1$ kept as an equation, not substituted into it). And
the constraint of the checking loop, $\al{\ldots}{k \mapsto k \le a'[k]}$, has been abstracted by
\lemma{above\_iterations\_def} of the database into the atom \eqref{eq:rw4} before the read $a'[k]$ was
resolved, since no lemma resolved it: the fact is hidden from any later rule on $a'[k]$ (finding~3 of the
algorithms, the lemma that ``swallows'' $k < t$ constraints). In \texttt{rewnif}, the read was rewritten
first and the constraint resolved by itself.

\subsection{The diagnosis}
The module's lemma is \lemma{select\_current\_trace\_ahead} mirrored: the dependencies $\bl{\cnt(0, 1, n)}$,
$\bl{\cnt(s, -1, n)}$, the cases $\upd{1 \le i}{a[i + (-1) := F(i)][i := G(i)]}$ and $\upd{i < 1}{a[i := G(i)]}$,
and the right-hand side $G(j)$ for $s - \intf(n) < j \le s$, $F(j + 1)$ at $j = s - \intf(n)$ if $0 \le j$ and
$0 < n$, else $a_0[j]$. Four variants were run on the test with the lemma injected as a user file:
\begin{center}\footnotesize
\begin{tabular}{@{}llll@{}}
\toprule
variant & step of the counter & position of the store ahead & result \\
\midrule
as in the module & literal $-1$ & $i + (-1)$ & does not fire \\
a & literal $-1$ & $i - 1$ & does not fire \\
b & variable $d$, condition $d = -1$ & $i + (-1)$ & does not fire \\
e & literal $-1$ & $P(i)$, pinned by $\forall i.\ P(i) = i - 1$ & does not fire \\
d & variable $d$, condition $d = -1$ & $P(i)$, pinned & \textbf{fires twice, closed} \\
\bottomrule
\end{tabular}
\end{center}
Both changes are needed, and both are patterns of Section~\ref{sec:principles:matcher} that the algorithms
had established: a literal $-1$ in a counter pattern is not matched (principle~8 of \texttt{NOTES.md},
\texttt{reverse.dfy}, whose counter down has its step as a variable with $d = -1$ as a side condition), and
arithmetic on a bound variable in a left-hand side, here the position $i - 1$ of the store, is not matched
and is pinned (Section~\ref{sec:merge}, the swaps at $j - 1$ of the insertion sort). The summary had excluded
the literal $-1$ because the database matches $\cnt(i, -1, n)$ elsewhere; it does so as a term, not as a
dependency pattern matched through the context of the goal, which is where it stands here. With variant d the
test closes with $s := \intf(n) + i'$, $F := \lambda z.\ z - 2$, $G := \lambda z.\ z$, also under
\texttt{min}; a false assertion \texttt{a[i] >= i + 1} leaves a goal. The variant does not cover
\texttt{rewnifrev2}, which stores ahead \emph{upwards} ($a[i + 1]$) in a loop that runs downwards, a third
shape, whose cell $j$ holds the value stored ahead by the iteration below it.

\subsection{The idea}
The database lemma is right and the engine applies it to the upward loop; the downward loop fails on two
literal patterns whose workarounds the algorithms had found in another guise. The lesson for the database is
that lemma statements must follow the matcher's patterns as much as the goal's forms: steps of counters as
variables with side conditions, positions computed from the bound variable as pinned higher-order
variables. The lesson for the engine is the same as finding~12 and finding~34 of the algorithms: a match that
fails without a visible reason costs more than its cause deserves, and a flag that reports which part of a
left-hand side did not match would have answered this in one run. The restated lemma is a candidate for the
module \texttt{strides}, which this document does not change.
